\section*{第 \ref{ch:g4:conductors-insulators} 章——导体与绝缘体}

\begin{solution}{exo:g4:conductors-insulators:1}
\omterm{def:g4:conductors-insulators:conductor}{导体}让电流通过它，\omterm{def:g4:conductors-insulators:insulator}{绝缘体}则把它挡住。
\omterm{def:g4:conductors-insulators:conductor}{导体}使\omterm{def:g3:first-electric-circuit:bulb}{灯泡}发光——它接通了回路。
\end{solution}

\begin{solution}{exo:g4:conductors-insulators:2}
\omterm{def:g4:conductors-insulators:conductor}{导体}：钢勺、锡纸、铜币。
\omterm{def:g4:conductors-insulators:insulator}{绝缘体}：软木塞、玻璃、塑料砖。
\end{solution}

\begin{solution}{exo:g4:conductors-insulators:3}
让两根导线的自由端直接接触：\omterm{def:g3:first-electric-circuit:bulb}{灯泡}必须亮起。
若测试器本身有故障，所有物体看起来都会是\omterm{def:g4:conductors-insulators:insulator}{绝缘体}，
这一检查可以避免错怪无辜的物体。
\end{solution}

\begin{solution}{exo:g4:conductors-insulators:4}
要么该物体是\omterm{def:g4:conductors-insulators:insulator}{绝缘体}——要么是测试操作草率：
导线端碰到的是油漆、清漆或污垢，而没有紧压在物体的
裸露材质上（或者连接松了）。
\end{solution}

\begin{solution}{exo:g4:conductors-insulators:5}
铜导电：它是电流通往\omterm{def:g3:first-electric-circuit:bulb}{灯泡}的通道。
塑料绝缘：它把电流留在导线内，
让手指安全地远离电流。
\end{solution}

\begin{solution}{exo:g4:conductors-insulators:6}
小金属片导电——\omterm{ex:g3:first-electric-circuit:switch}{开关}闭合时，它必须接通回路。
外壳绝缘——手指整天都要按压它。
没有金属片就不会发光；没有外壳就没有安全。
\end{solution}

\begin{solution}{exo:g4:conductors-insulators:7}
是的——铝是金属，而所有金属都导电。注意：
被\omterm{def:g1:magnets:magnet}{磁铁}\omterm{def:g1:magnets:magnet}{吸引}与导电是两回事，
\omterm{def:g1:magnets:magnet}{磁铁}“只吸铁”的规律，并不能说明物体是否导电。
\end{solution}

\begin{solution}{exo:g4:conductors-insulators:8}
木材（和油漆）是\omterm{def:g4:conductors-insulators:insulator}{绝缘体}，但铅笔芯里的石墨是\omterm{def:g4:conductors-insulators:conductor}{导体}——
一种少见的非金属\omterm{def:g4:conductors-insulators:conductor}{导体}，
贯穿整支铅笔，从笔尖一直延伸到笔尾。
\end{solution}

\begin{solution}{exo:g4:conductors-insulators:9}
\omterm{def:g4:conductors-insulators:insulator}{绝缘体}在保护人：若工具的金属部分
碰到了带电的东西，较厚的塑料手柄
可以阻止电流传到手上。
\end{solution}

\begin{solution}{exo:g4:conductors-insulators:10}
电流必须走完整条通路。链条中，金属环与金属环相连，
形成一条完整的金属通路。串珠项链中，相邻两颗金属珠
之间隔着一段棉线结——即\omterm{def:g4:conductors-insulators:insulator}{绝缘体}——因此每颗珠子处
通路都被切断：尽管有那么多\omterm{def:g4:conductors-insulators:conductor}{导体}，
灯也是暗的。
\end{solution}

\begin{solution}{exo:g4:conductors-insulators:11}
潮湿的皮肤比干燥的皮肤导电好得多，所以湿手
为电流提供了一条容易穿过身体的通路。\omterm{def:g3:first-electric-circuit:battery}{电池}测试器是安全的，
因为\omterm{def:g3:first-electric-circuit:battery}{电池}的推力很小；而墙上插座的推力
要大数百倍——通过潮湿的皮肤，它可能致命，
所以规则允许玩\omterm{def:g3:first-electric-circuit:battery}{电池}，却一直禁止接触插座，
始终如此。
\end{solution}
